\section{Introduction}
\label{sec:intro}

The statistical mechanics of learning has produced a remarkable body of exact
results: the storage capacity of the perceptron, the replica solution of
generalised linear estimation, the online dynamics of the soft committee
machine, the asymptotics of random-feature and kernel regression. What these
results share is not a technique but a \emph{regime}. They live in the
proportional high-dimensional limit, with i.i.d.\ Gaussian inputs, shallow
architectures, and a loss whose minimiser can be characterised. Almost every
system that motivates the field today lies outside that regime.

The usual response is to wait for theory to arrive, and meanwhile to describe
modern phenomena --- grokking, double descent, emergence, collapse --- with
statistical-mechanics vocabulary borrowed on credit. Words like ``phase,''
``order parameter,'' and ``critical point'' are applied to curves that have
never been subjected to the finite-size analysis that gives those words
content. The result is a literature rich in suggestive plots and poor in
falsifiable statements.

This paper takes the opposite route. If the analytic calculation is not
available, the finite-size numerical experiment that a theory would ultimately
have to reproduce still is, and it can be done to the standards of
computational statistical mechanics rather than those of a learning-curve
figure. We build an \emph{atlas}: a broad grid of learning settings, each
measured with registered order parameters on a control grid fine enough to
resolve a boundary, across a scale path long enough to extrapolate, over a
disorder ensemble large enough to give the fluctuations meaning.

\paragraph{What makes this possible now is not compute but design.}
The measurements here are deliberately small --- widths of tens to hundreds,
models of thousands of parameters --- because the quantities that decide
whether a transition exists are not accuracy on a large model but the
\emph{scaling of fluctuations} across many independent small ones. The binding
constraint in this kind of study is the number of independent runs, not the
size of any one of them. That inverts the usual resource allocation: this
release trains \AtlasTotalRuns{} models across \AtlasTotalConditions{}
conditions on a laptop, where a single conventional experiment of comparable
cost would have produced one learning curve.

\paragraph{Anchors first.}
A measurement pipeline that cannot reproduce a known result should not be
trusted on an unknown one. The atlas therefore opens with settings whose
answers are known exactly --- the Gardner storage capacity and the
Gaussian-mixture error formula --- and reports the measured quantity against
the closed form, with the residual as a first-class number. Only then does it
move outward, through settings with partial theory, to settings with none.

\paragraph{Negative and unresolved outcomes are results.}
The decision procedure in \S\ref{sec:gates} can return ``no resolved
response'' and ``critical response, boundary unresolved'' as readily as
``transition supported,'' and this release does return all of them. One case is
worth stating in the introduction because it shaped the method: our flagship
family initially produced a confident transition claim that was an artefact of
two design choices --- a control grid too coarse to place a peak, and a fixed
epoch budget under which wider models were simply optimised less. Fixing the
first turned a censored boundary into a resolvable one; fixing the second
removed a spurious size dependence. The consistency requirement that caught the
artefact (\S\ref{sec:gates}) is now a gate every family must pass.

\paragraph{Contributions.}
\begin{enumerate}
  \item A portable set of function-space order parameters
    (\S\ref{sec:observables}) that can be evaluated on any differentiable
    learner, together with the exact identity \eqref{eq:epsg} used as a
    continuous unit test on the measurement chain.
  \item A decision procedure (\S\ref{sec:gates}) that separates
    \emph{sharpening} from \emph{location}, requires independent finite-size
    signatures to agree before they accumulate, and names the outcome when they
    do not.
  \item An empirical atlas of \AtlasTotalFamilies{} experiment families across
    \AtlasTotalBlocks{} blocks --- \AtlasAnchorFamilies{} anchored by exact
    theory and \AtlasNoTheoryFamilies{} with no theory at all --- each with a
    fine control grid, a multi-size scale path, and \AtlasSpecializationSeeds{}
    outer seeds per condition.
  \item Quantitative targets for future theory: measured exponents, boundary
    locations, and scaling paths, each reported with the interval and the
    experiment that would falsify it.
\end{enumerate}

The claim of this paper is not that these systems have been solved. It is that
the question ``does this system have a phase boundary, and where?'' can be
asked well enough to be answered, or well enough to be shown unanswerable with
present evidence --- and that doing so is worth more than another suggestive
curve.
